\section{Post-Hoc Leakage Audit and Follow-Up Verification}
\label{sec:leakage-audit}

Sections~\ref{sec:results}.1--4 report the original XAI battery, run on a split believed at the time to be free of near-duplicate leakage. A subsequent, independent audit found that belief incomplete. This section tells that story in three parts. Section~\ref{sec:stage2} reports the correction itself: a cluster-aware re-split, a genuine held-out test set, and the corrected accuracy this yields. The methodological questions that correction raised are then answered in Section~\ref{sec:followup-audits}, in the order they arose: is the batch/session inference this whole audit rests on actually justified, is the clustering threshold that defines it well-chosen, does the batch-shortcut finding survive the same leakage fix, and how far does classifier accuracy itself collapse under noise. Finally, Section~\ref{sec:stage2-full-battery} closes the gap those first two parts leave open, re-running the complete five-method battery, not just the one method each earlier check re-verified, so every number in the paper's XAI comparison is finally about the same, honestly-evaluated classifier its accuracy claims are about. This is presented as its own section rather than folded into Section~\ref{sec:results} because it changes the paper's headline accuracy figure and because the story is easier to follow as a self-contained unit than scattered through an already-long results section.

\subsection{Stage 2: cluster-aware split with a genuine held-out test set}
\label{sec:stage2}

The results in Section~\ref{sec:results} come from the Stage~1 pipeline: images de-duplicated by single-orientation perceptual hashing (Section~\ref{sec:methods}.1), then split 70/30 into training and validation, with every number in that section reported on that validation set. A subsequent, independent audit found this correction incomplete in two ways. First, 96.73\% of images share a filename-prefix ``session'' that the Stage~1 split did not keep together, so near-identical frames from the same imaging session could land on both sides of the train/validation boundary. Second, the dataset's own source publication confirms images were released after augmentation by rotation, flipping, and brightness/contrast/saturation jitter, and a rotation-aware (dihedral) hash re-check found 46\% of the corpus (4,918 of 10,692 images) has at least one cross-split duplicate that single-orientation hashing could not detect.

A Stage~2 pipeline addresses both gaps together. Every image is assigned to a duplicate cluster using an 8-orientation (dihedral) perceptual hash (1,570 clusters from the full 10,692-image corpus), and the corpus is re-split at the cluster level, 70/20/10, so no cluster crosses a train/validation/test boundary and, unlike Stage~1, a test set is held out and touched exactly once, after every other design decision (epoch budget, early-stopping criterion) is already frozen. This yields 7,470 training, 2,138 validation, and 1,084 test images. A primary model (seed~42) and three replicate models (seeds~7, 123, 2027) are trained on the identical 7,470-image training set, varying only the random seed and data-loader shuffling order---correcting the Stage~1 replicate design, in which the three replicates were trained on disjoint subsets of 2,749, 2,089, and 2,661 images respectively and therefore could not isolate seed-to-seed reproducibility from data-composition effects (Limitation~2b). Training uses up to 20 epochs with early stopping on validation macro-F1 (patience~5), rather than Stage~1's fixed 10-epoch budget.

\begin{table}[!ht]
\centering
\caption{Stage 2: Held-Out Test Performance, Cluster-Aware 70/20/10 Split (Test $n=1084$). All four models see identical training data; only the random seed varies.}
\label{tbl:stage2-test}
\begin{tabular}{lccc}
\toprule
Model & Test Accuracy & Macro-F1 & Epochs to best \\
\midrule
Primary (seed~42) & 0.7030 & 0.7039 & --- \\
Replicate (seed~7) & 0.7020 & 0.6992 & --- \\
Replicate (seed~123) & 0.7482 & 0.7486 & --- \\
Replicate (seed~2027) & 0.7491 & 0.7438 & --- \\
\midrule
\textbf{Mean~$\pm$~Std} & \textbf{0.7256~$\pm$~0.0235} & \textbf{0.7239~$\pm$~0.0221} & \\
\bottomrule
\end{tabular}
\end{table}

Table~\ref{tbl:stage2-test} reports the real, unbiased test performance this design supports: a four-model mean accuracy of 72.56\%~$\pm$~2.35\% and macro-F1 of 72.39\%~$\pm$~2.21\%, all four seeds agreeing within roughly a five-point band. This is substantially lower than the Stage~1 headline of 93.05\% accuracy, and it is the more trustworthy of the two numbers precisely because Stage~1's validation set was not fully protected from cluster-level leakage. It is also markedly higher than the 29.59\% single-run figure produced by an earlier, more aggressive session-level-only correction (Section~\ref{sec:limitations}), which re-split the existing Stage~1 train/validation files at the filename-session level without also addressing the rotation-augmented-duplicate gap or retraining with a longer schedule; that figure should be read as a lower bound from an under-specified re-split rather than the pipeline's ceiling. The Stage~2 result in Table~\ref{tbl:stage2-test} is the paper's best current estimate of what this classifier, honestly evaluated, actually achieves on nine-class gallbladder ultrasound classification.

\textit{XAI Finding (9): the null-control test for the consistency metric.} A structural concern raised against the Stage~1 consistency scores in Table~\ref{tbl:5} is that Grad-CAM-family heatmaps end in a ReLU and are then min-max normalised to $[0,1]$, so two \textit{unrelated} non-negative heatmaps of the same shape may already have high cosine similarity as a mathematical artefact, independent of anything the model has learned. Two null controls quantify this floor directly. First, cross-model consistency is measured between four \textit{untrained} ResNet-50 models (ImageNet backbone, randomly initialised nine-class head, no fine-tuning): mean cosine similarity $0.333\pm0.231$ over 240 pairwise comparisons across 40 stratified test images. Second, cosine similarity is measured between pairs of purely synthetic heatmaps with no model, image, or learning involved: $0.755\pm0.044$ for independent uniform-random $[0,1]$ maps, and $0.325\pm0.107$ for ReLU-thresholded standard-Gaussian noise, min-max normalised the same way Grad-CAM's own output is---the more representative control, since it reproduces Grad-CAM's actual non-linearity rather than assuming uniform noise. Against this pair of floors ($0.33$ untrained-model, $0.33$ ReLU-Gaussian-synthetic), the Stage~2 same-data cross-replicate consistency of $0.848$, $0.816$, and $0.863$ (mean $\approx0.843$, per-replicate-mean std $=0.0196$; Section~\ref{sec:results}.3's Grad-CAM++ score, recomputed here under a corrected same-data replicate design) clears both floors by roughly a factor of 2.5, rather than sitting close to them. The consistency claim therefore survives the null-control test the metric had not previously been checked against. Two caveats qualify that conclusion: this recomputation used Grad-CAM specifically as the representative CAM-family method, not the full five-method comparison of Table~\ref{tbl:5}, and the uniform-random floor of $0.755$, higher than either the untrained-model or ReLU-Gaussian floor, shows the choice of null distribution itself matters and should be reported alongside any future consistency claim rather than assumed.

The Stage~2 per-replicate-mean standard deviation reported above ($0.0196$) also corrects a labelling error in the Stage~1 tables. Table~\ref{tbl:5}'s std column was computed by pooling every individual image-level score across all three replicates and taking the standard deviation of that pooled set, not the standard deviation of the three replicate-level \textit{means} the table and its caption describe. Recomputed the way the caption implies, the three Stage~2 replicate means $\{0.848, 0.816, 0.863\}$ give a std of $0.0196$; the pooled-image-level std over the same data is $0.186$, an order of magnitude larger. Both quantities are legitimate to report, but they answer different questions---how much do the three replicates' \textit{average} agreement levels differ from one another, versus how much does agreement vary image-to-image within a replicate---and only the first is a std-of-means in the sense Table~\ref{tbl:5}'s caption claims.

\subsection{Follow-up audits: threshold sensitivity, batch-shortcut re-verification, and noise robustness}
\label{sec:followup-audits}

The Stage~2 correction itself raised four further questions, each targeted at a different assumption the correction depends on, plus one question raised separately by the stability results in Section~\ref{sec:results}.4. Two concern the cluster-aware split's own foundations: is the filename-prefix inference it is built on actually justified, and is its clustering threshold well-chosen. Two concern whether the batch-shortcut finding of Section~\ref{sec:results}.4 survives the same leakage fix, and whether it reaches the disease classifier's own features rather than being confined to a separately-trained proxy. The fifth asks how far classification accuracy itself, not just explanation stability, degrades under noise. All five are addressed here using the same dihedral cluster manifest and Stage~2 primary checkpoint as Section~\ref{sec:stage2}.

\textbf{Is the filename-prefix inference justified?} Section~\ref{sec:methods}.1 infers that each filename prefix denotes one contiguous acquisition or export session, an inference previously argued only from filename structure: prefixes follow a consistent pattern, indices are contiguous within a prefix, and no prefix spans two disease-class folders. A reviewer correctly noted this argument is partly circular, since prefixes assigned per class folder at export time would trivially satisfy the single-class property regardless of whether they reflect genuine acquisition sessions. Every one of the 10,692 released images carries embedded EXIF metadata (all exported via ACDSee~Pro~6, per the embedded software tag), including a DateTimeOriginal timestamp, which is direct, non-inferred evidence independent of the filename pattern itself. Grouping images by filename prefix and examining each group's timestamp range shows a median within-prefix capture-time span of just 5.7 minutes across 220 prefixes with usable EXIF data, consistent with one short imaging or export session. More decisively: comparing every pair of prefixes that occur on the same calendar day for whether their time windows overlap at all finds \textbf{zero overlapping pairs out of 3,312 same-day prefix pairs checked}. No two filename-prefix groups ever share a moment in time. An arbitrary per-class export label would be expected to produce interleaved or overlapping time windows across prefixes; this instead is direct evidence that each prefix corresponds to its own distinct, temporally exclusive session. Thirteen of the 220 prefixes (5.9\%) run longer than an hour, three substantially so (5.8, 7.1, and 7.7 hours); we report these rather than exclude them, since they may reflect a genuine pause, batched review-and-export, or some other export-side cause this metadata alone cannot distinguish. This does not replace contacting the original providers for their actual export convention, which remains unresolved, but it is direct, quantitative, non-circular support for an inference the rest of this audit depends on.

\textbf{Is the clustering threshold that split rests on well-chosen?} Section~\ref{sec:methods}.1's clustering uses a Hamming-distance threshold of 4 bits, tightened from an initial choice of 8 after threshold~8 produced a single-linkage chaining artefact: a 498-image cluster spanning multiple disease classes, later traced to a chain of transitively-connected near-threshold matches rather than genuine duplicates. Table~\ref{tbl:phash-sweep} sweeps this threshold from 4 to 12 bits on the same cached dihedral hashes, so the comparison isolates the threshold's effect alone. Two numbers move together as the threshold loosens: cluster count roughly triples by threshold~8 (1,570~$\rightarrow$~2,252) and explodes by threshold~10 (5,121 clusters, median size dropping back to 1, meaning most images become effective singletons again), while mixed-class clusters, a proxy for false-positive merges, rise from 4 at threshold~4 to 26 at threshold~8. Threshold~4 is the more conservative choice on both counts: smallest, most internally-consistent clusters, fewest cross-class merges, at the cost of a somewhat larger effective sample size than a looser threshold would give. What this sweep does not do is retrain a classifier at each threshold to report Stage~2 accuracy per threshold, which would require up to five full training runs; the cluster-count and mixed-class sensitivity reported here is a lower-cost but real partial answer, not a substitute for that larger check.

\begin{table}[!ht]
\centering
\caption{pHash Clustering Threshold Sensitivity (Hamming distance, dihedral 8-orientation hashing, 10,692 raw images). Threshold 4 is used throughout this paper; thresholds 6--12 are shown for comparison only.}
\label{tbl:phash-sweep}
\begin{tabular}{lccc}
\toprule
Threshold (bits) & Final Clusters & Mixed-Class Clusters & Median Cluster Size \\
\midrule
4 (used here) & 1,570 & 4 & 6 \\
6 & 1,573 & 14 & 6 \\
8 & 2,252 & 26 & 6 \\
10 & 5,121 & 32 & 1 \\
12 & 9,299 & 12 & 1 \\
\bottomrule
\end{tabular}
\end{table}

With the split itself examined, the remaining three questions turn to the batch-correlated shortcut Section~\ref{sec:results}.4 first raised. \textbf{Does the batch-ID diagnostic survive the same leakage fix that motivated Stage~2?} That diagnostic, a $K$-way classifier trained to predict which acquisition batch a Carcinoma image came from, all images sharing one disease label, originally used a random image-level split within each batch. A review pointed out this could let near-duplicate frames from the same acquisition sit on both sides of that split, inflating its reported 100\% accuracy. Re-running the same 10-batch diagnostic with the dihedral cluster manifest applied, so no duplicate cluster crosses the train/validation boundary (608 training, 292 validation images), gives \textbf{82.88\% accuracy}, 8.3$\times$ the 10\% chance level and well down from the original 100\%, but still far above chance. The review's concern was partly correct: some of the original 100\% figure was inflated by near-duplicate-frame memorisation. Its underlying conclusion also survives: a substantial, learnable, non-anatomical batch signature remains even under a leakage-controlled split.

\textbf{Does that signature generalise beyond memorised batches, and does it reach the disease classifier's own features?} Two further tests answer this directly, since the paper's only other shortcut audit (Section~\ref{sec:results}.4) tests exclusively for margin/border artefacts and cannot, by its own design, detect a batch-correlated one. First, a 6-way batch classifier trained on a subset of batches reaches 90.31\% held-out accuracy on unseen images from those same batches (chance 16.67\%). Its mean top-1 confidence on four entirely unseen batches (73.45\%) stays substantial relative to its confidence on seen-batch held-out images (89.77\%), suggesting the signature partially generalises rather than being purely memorised per batch. Second, and more directly relevant to whether the \textit{disease} classifier itself is affected: a linear probe on the Stage~2 primary disease classifier's own \texttt{layer4} features, the same features that classifier uses to make its actual nine-class decision, recovers batch identity with \textbf{74.75\% accuracy} against 6 batches (chance 16.67\%, 4.5$\times$ chance). This is direct evidence that the disease classifier's learned representation itself correlates with batch identity, not merely that a separately-trained proxy classifier is capable of learning it. The shortcut concern therefore extends to the classifier this paper's own XAI battery is evaluating explanations for.

The fifth question is unrelated to the split or the batch shortcut, and comes instead from Section~\ref{sec:results}.4's stability test, which reported that Gaussian noise at $\sigma=0.03$ changes the predicted class on 29.2\% of images under the Stage~1 model but stopped there: no direct accuracy number, no sweep beyond that one noise level. \textbf{How far does classification accuracy itself collapse under noise?} Table~\ref{tbl:noise-sweep} answers this directly, sweeping accuracy and prediction-change rate across seven noise levels on the Stage~2 primary checkpoint. Accuracy on clean images is 90.0\% on this 180-image stratified subset. It falls to 68.3\% by $\sigma=0.03$ and to 20.6\%, barely above the 11.1\% chance level, by $\sigma=0.12$, a noise level that remains visually subtle. The fraction of predictions that flip relative to the clean prediction tracks this drop closely, 7.2\% at $\sigma=0.01$ rising to 82.2\% at $\sigma=0.12$, confirming the accuracy collapse and the prediction instability are the same underlying phenomenon rather than two independent effects. This is a substantially more severe robustness finding than anything in the XAI comparison itself: as Section~\ref{sec:results}.4 already notes, classifier-level noise robustness, not explanation-level stability, is the more urgent target for future model development.

\begin{table}[!ht]
\centering
\caption{Noise-Robustness Sweep, Stage 2 Primary Checkpoint (180 stratified images, 20/class). Accuracy and prediction-change fraction both reported directly, closing a gap where only a single $\sigma$ was previously reported.}
\label{tbl:noise-sweep}
\begin{tabular}{lcc}
\toprule
Noise $\sigma$ & Accuracy & Pred.\ Changed from Clean \\
\midrule
0.00 (clean) & 0.900 & --- \\
0.01 & 0.878 & 0.072 \\
0.02 & 0.800 & 0.172 \\
0.03 & 0.683 & 0.328 \\
0.05 & 0.511 & 0.511 \\
0.08 & 0.333 & 0.683 \\
0.12 & 0.206 & 0.822 \\
\bottomrule
\end{tabular}
\end{table}

\subsection{The full five-method XAI battery, re-run on Stage 2}
\label{sec:stage2-full-battery}

Every result above re-verifies one method (Grad-CAM or Grad-CAM++) against the Stage~2 correction, not the complete five-method, six-analysis battery of Section~\ref{sec:results}. This leaves the paper's central comparison, which XAI method is most reliable, resting on the Stage~1 classifier alone, even though Section~\ref{sec:stage2} shows that classifier's accuracy was inflated by leakage. This subsection closes that gap: the full battery (consistency across all five methods, faithfulness, stability, the parameter-randomization sanity check, calibration, and the margin-shortcut audit) is re-run unchanged on the Stage~2 primary and replicate checkpoints and the materialized cluster-aware test set (Section~\ref{sec:stage2}), using the identical scripts that produced every Stage~1 number in Section~\ref{sec:results}, so the two sets of results are a fair comparison rather than a reimplementation. Table~\ref{tbl:stage2-battery} places every headline figure from both stages side by side.

\begin{table*}[!ht]
\centering
\caption{Full XAI Reliability Battery, Stage 1 (leakage-inflated split) vs. Stage 2 (cluster-aware, genuinely held-out test set). Stage 2 consistency and stability use the same 5-method / transform battery as Stage 1, on the materialized 1,084-image cluster-aware test set.}
\label{tbl:stage2-battery}
\begin{tabular}{llcc}
\toprule
Analysis & Metric & Stage 1 & Stage 2 \\
\midrule
Classifier & Test/validation accuracy & 0.9305 & 0.7256~$\pm$~0.0235 \\
\midrule
\multirow{5}{*}{Consistency $O^c$ (macro mean)}
 & Grad-CAM & 0.769 & 0.894 \\
 & Grad-CAM++ & \textbf{0.944} & 0.963 \\
 & Saliency Map & 0.681 & 0.707 \\
 & Eigen-CAM & 0.242 & 0.860 \\
 & Score-CAM & 0.937 & \textbf{0.964} \\
\midrule
\multirow{3}{*}{Faithfulness gap}
 & Grad-CAM & \textbf{0.480} & \textbf{0.487} \\
 & Grad-CAM++ & 0.473 & 0.475 \\
 & Score-CAM & 0.456 & 0.445 \\
\midrule
\multirow{3}{*}{Stability SSIM}
 & Brightness ($\times1.15$) & 0.995 & 0.994 \\
 & Translation (10~px) & 0.848 & 0.857 \\
 & Gaussian noise ($\sigma=0.03$) & 0.931 & 0.886 \\
\midrule
Sanity check & $|\rho|$ at layer4 (trained-model tap point) & 0.52 & 0.469 \\
\midrule
\multirow{2}{*}{Calibration}
 & ECE, before~$\rightarrow$~after temperature scaling & 0.034~$\rightarrow$~0.009 & 0.184~$\rightarrow$~0.081 \\
 & Brier, before~$\rightarrow$~after & 0.109~$\rightarrow$~0.106 & 0.511~$\rightarrow$~0.468 \\
\midrule
\multirow{2}{*}{Shortcut audit}
 & Accuracy, full~$\rightarrow$~cropped & 0.930~$\rightarrow$~0.712 & 0.703~$\rightarrow$~0.516 \\
 & Attribution mass in margin ring & 0.218 & 0.239 \\
\bottomrule
\end{tabular}
\end{table*}

Two patterns emerge, and they point in opposite directions for how much the Stage~1 XAI comparison can be trusted.

First, the \textit{relative ranking} across methods is almost unchanged. Grad-CAM++ and Score-CAM remain the two most cross-replicate-consistent methods by a wide margin (0.963 and 0.964 vs. Grad-CAM's 0.894, Saliency's 0.707); Grad-CAM still numerically edges out the other two CAM methods on faithfulness gap (0.487 vs. 0.475 and 0.445); and the sanity check still shows the expected sharp similarity collapse once \texttt{layer4} is randomized ($|\rho|$ falling from 0.921 at \texttt{fc} to 0.469 at \texttt{layer4}), the same qualitative pattern Fig.~\ref{fig:12} shows for Stage~1. Every one of the paper's method-comparison conclusions in Section~\ref{sec:results}.5 therefore survives the leakage correction unchanged in direction, even though the underlying classifier's accuracy does not.

Second, the \textit{absolute} numbers move, and not uniformly. Consistency scores rise slightly for every method under Stage~2, most strikingly Eigen-CAM (0.242~$\rightarrow$~0.860). This is a re-verification of the caveat already stated in Table~\ref{tbl:5}'s footnote rather than a new finding: Eigen-CAM's low Stage~1 score reflects a class-agnostic method being scored by a class-conditional metric, and a different, harder-to-control random split naturally perturbs that mismatch's severity. Faithfulness gaps and stability SSIMs move by only a few points in either direction, consistent with these being properties of the CAM methods' mechanics more than of the specific classifier they are applied to. Calibration is the one analysis where Stage~2's absolute numbers are substantially worse, not just different: ECE before scaling roughly quintuples (0.034~$\rightarrow$~0.184). This is the expected consequence of accuracy falling from 93.05\% to 70.30\% on the primary model, since a less accurate classifier has more room to be over-confident about its wrong answers; temperature scaling still recovers a comparable relative improvement (78\% ECE reduction under Stage~2 vs. 74\% under Stage~1). The shortcut audit's crop-ablation drop is smaller in absolute terms under Stage~2 (0.703~$\rightarrow$~0.516, a further 18.7-point drop) than Stage~1 (0.930~$\rightarrow$~0.712, a 21.8-point drop), but represents a larger \textit{relative} collapse (26.6\% vs. 23.5\% relative accuracy loss). The margin-shortcut-vs-genuine-anatomy ambiguity already flagged in XAI Finding~(8) is, if anything, slightly sharper under the corrected split, not resolved by it.

This re-run supports a practical conclusion that could not have been stated with confidence before it: the paper's central claim, that Grad-CAM++ and Score-CAM are the more reliable explanation methods on the consistency axis, while all three CAM methods are comparably faithful and none dominates on every axis, is not an artefact of the Stage~1 leakage. It holds on the same classifier's honestly-evaluated, genuinely held-out counterpart. What does \textit{not} survive unchanged is any claim resting on the classifier's absolute accuracy or confidence, which Section~\ref{sec:stage2} already corrects, and calibration quality, which degrades together with accuracy as expected rather than independently.

Two scope notes apply to this full re-run. It uses the same stratified sample sizes as Section~\ref{sec:results} (72 images for consistency, 144 for faithfulness, 216 for stability and the shortcut audit, 108 for the sanity check, a stratified half-split of the 1,084-image test set for calibration), drawn from the materialized Stage~2 test set rather than the full corpus, for the same computational reasons given in Section~\ref{sec:methods}.3. And it does not re-run the failure-case visual analysis of Section~\ref{sec:results}.2, which would require a fresh selection of high-confidence Stage~2 misclassifications and is left to future work (Section~\ref{sec:future-research}).
